\section{Conclusion}
% 我们提出了MAP-STAR，一个能够学习可迁移图结构的新颖的区域表示学习框架。与传统的在固定图结构上进行表示学习的方法不同，我们的框架提出在异质采样子图上进行消息传递与表示学习，从而克服传统表示学习方法无法进行迁移学习的困难。为了使得模型在子图的学习上捕获更一般的图结构模式信息，我们设计了Multi-view Pattern-bank Subgraph Encoder，首先通过节点间的异质关系聚合不同邻居节点的表征，之后基于多视图的Pattern Bank进行表征聚合与参数更新，从而在模型中注入包含区域邻接结构、人口流动模式、区域属性结构等多维度信息。Extensive experiments说明了我们提出方法能够在不同城市的不同下游任务中取得较好的迁移学习效果，体现我们的方法在区域迁移学习、乃至更一般的图结构迁移学习的巨大潜力。我们工作未来可能的研究方向包括探究不同下游任务对不同子图采样策略的attention大小，探究方法框架在更一般graph learning任务上的效果等。

We propose MAP-STAR, a novel region representation learning framework designed to acquire transferable graph structures across cities. Unlike conventional methods that learn node embeddings on a fixed full graph, MAP-STAR shifts representation learning to heterogeneous sampled subgraphs, which alleviates structural entanglement and improves robustness under cross-city distribution shifts. By leveraging a multi-view pattern-bank encoder, MAP-STAR aggregates complementary signals from geographic adjacency, human mobility, and regional attributes, distilling reusable structural patterns into compact and semantically expressive representations. Extensive experiments on multiple cities and diverse downstream tasks demonstrate that MAP-STAR consistently achieves superior performance in both single-city and cross-city settings, highlighting its effectiveness for regional representation learning and, more broadly, graph-structure transfer in heterogeneous urban environments. Future work will explore task-adaptive subgraph sampling strategies and extend the framework to more general graph transfer learning scenarios beyond urban computing.

